\begin{lemma}[Finite dependence on the enumeration orbit]
For fixed $r,n,k$, the copied value
$\mathsf{PowerCopiedLeft}_r(h,x,n,k)$ depends only on the finitely many
values of $h$ on the first $n+k+2$ shifts of $x$.  Thus, if $x$ and $y$
give the same $h$-values at all these shifts, then
\[
\mathsf{PowerCopiedLeft}_r(h,x,n,k)
=\mathsf{PowerCopiedLeft}_r(h,y,n,k).
\]
\end{lemma}
\begin{proof}
Induct on $k$.  At depth zero, the definition in
\noderef{PowerCopiedLeft} uses only the values at shifts $n$ and $n+1$,
which are covered by the hypothesis.  At depth $k+1$, the two recursive
inputs are the values at $(n,k)$ and $(n+1,k)$.  The induction hypotheses
apply to these inputs because their required bounds are respectively
$n+k+2$ and $n+k+3$, exactly the bound
$n+(k+1)+2$.  The same deterministic response function from
\noderef{PowerIResponse} is applied to equal recursive inputs, so the
outputs are equal.
\end{proof}
